\documentclass[a4paper,12pt]{article}
\usepackage[T2A]{fontenc}
\usepackage[utf8]{inputenc}
\usepackage[russian]{babel}
\usepackage{amsmath,amssymb,mathtools}
\usepackage{DejaVuSans}
\renewcommand{\familydefault}{\sfdefault}
\usepackage{icomma}
\usepackage[a4paper,margin=17mm,headheight=25pt,headsep=7mm,footskip=11mm]{geometry}
\usepackage{microtype}
\usepackage{needspace}
\usepackage{tikz}
\usetikzlibrary{calc,arrows.meta,positioning,shapes.geometric}
\usepackage{tabularx,booktabs,array}
\usepackage{enumitem}
\usepackage{xcolor}
\usepackage{fancyhdr}
\usepackage{setspace}
\usepackage{hyperref}
\usepackage{parskip}

\definecolor{accent}{HTML}{2D6C69}
\definecolor{darkaccent}{HTML}{174A4B}
\definecolor{textgray}{HTML}{252B2D}
\definecolor{softgray}{HTML}{687477}
\definecolor{linegray}{HTML}{CBD8D7}
\hypersetup{colorlinks=true,linkcolor=darkaccent,urlcolor=darkaccent}
\setlength{\parindent}{0pt}
\setstretch{1.07}
\setlist[itemize]{leftmargin=1.6em,itemsep=0.18em,topsep=0.18em}
\setlist[enumerate]{leftmargin=1.65em,itemsep=0.45em,topsep=0.30em}
\pagestyle{fancy}
\fancyhf{}
\fancyhead[L]{\small\color{softgray} Умножение дробей и задачи}
\fancyhead[R]{\small\color{softgray} 20.10.26}
\fancyfoot[C]{\small\color{softgray}\thepage}
\renewcommand{\headrulewidth}{0.3pt}
\renewcommand{\headrule}{\hbox to\headwidth{\color{linegray}\leaders\hrule height \headrulewidth\hfill}}
\newcommand{\mixed}[3]{#1\,\frac{#2}{#3}}
\newcommand{\sectiontitle}[1]{%
\Needspace{5\baselineskip}%
\vspace{0.34em}{\large\bfseries\color{darkaccent}#1}\par
\vspace{0.16em}{\color{linegray}\hrule}\vspace{0.33em}}
\newcommand{\rulelabel}[1]{\textbf{#1.}}
\newcommand{\tip}[1]{\textbf{Обратите внимание.} #1}
\newcommand{\keyidea}[1]{\textbf{Главная мысль.} #1}

\begin{document}
\color{textgray}
\setlength{\abovedisplayskip}{5pt}
\setlength{\belowdisplayskip}{5pt}

\hypersetup{pageanchor=false}
\begin{titlepage}
\thispagestyle{empty}
\centering
\vspace*{26mm}
{\small\color{softgray} МАТЕМАТИКА \quad • \quad ПОВТОРЕНИЕ И ПРАКТИКА}\par
\vspace{12mm}
{\Huge\bfseries\color{darkaccent} Чек-лист}\par
\vspace{9mm}
{\LARGE\bfseries Умножение дробей\\[2mm] и текстовые задачи}\par
\vspace{5mm}
{\large Доли, проценты, движение, распределительное свойство}\par
\vspace{18mm}
{\large 20.10.26}\par
\vspace{9mm}
{\normalsize Тренировочный блок: Путь А · 10 заданий}\par
\vfill
{\large\bfseries Лёвин Артём Александрович}\par
\vspace{2mm}
{\normalsize эксклюзивно для Нади Клименко}\par
\vspace{20mm}
\begin{tikzpicture}[remember picture,overlay]
\draw[accent,line width=1.0pt]
($(current page.south west)+(18mm,20mm)$)
.. controls ($(current page.south west)+(64mm,30mm)$)
and ($(current page.south east)+(-68mm,11mm)$)
.. ($(current page.south east)+(-18mm,20mm)$);
\end{tikzpicture}
\end{titlepage}

\hypersetup{pageanchor=true}
\setcounter{page}{1}
\sectiontitle{1. Умножение обыкновенных дробей}
\keyidea{Числители перемножаем между собой, знаменатели --- между собой. Если возможно, сокращаем множители \emph{до} умножения.}
\[
\frac ab\cdot\frac cd=\frac{ac}{bd},
\qquad b\ne0,\quad d\ne0.
\]
\rulelabel{Пример} Сначала сократим 2 и 4:
\[
\frac23\cdot\frac14=\frac{2\cdot1}{3\cdot4}
=\frac{1\cdot1}{3\cdot2}
=\boxed{\frac16}.
\]
Точно так же можно перемножать три и более дроби:
\[
\frac23\cdot\frac94\cdot\frac18
=\frac{2\cdot9\cdot1}{3\cdot4\cdot8}
=\boxed{\frac3{16}}.
\]
\tip{Сокращать можно числитель одной дроби со знаменателем \emph{другой}. Сокращать два числителя между собой нельзя.}

\sectiontitle{2. Умножение на целое число и смешанные числа}
Целое число представляем дробью со знаменателем 1:
\[
\frac23\cdot9=\frac23\cdot\frac91=\boxed6.
\]
Смешанное число сначала переводим в неправильную дробь:
\[
\mixed{2}{1}{3}=\frac{2\cdot3+1}{3}=\frac73,
\qquad
\mixed{2}{1}{3}\cdot\frac92
=\frac73\cdot\frac92=\frac{21}{2}
=\boxed{\mixed{10}{1}{2}}.
\]
\tip{В конце выделите целую часть, если получилась неправильная дробь и этого требует форма ответа.}

\sectiontitle{3. Степень дроби и значение скобок}
Если в квадрат возводится \emph{вся дробь}, скобки обязательны:
\[
\left(\frac12\right)^2=\frac12\cdot\frac12=\boxed{\frac14},
\qquad
\frac{1^2}{2}=\boxed{\frac12}.
\]
Это разные выражения. В общем случае:
\[
\left(\frac ab\right)^2=\frac{a^2}{b^2},\quad b\ne0.
\]
\tip{Степень относится только к тому выражению, которое непосредственно ей предшествует.}

\sectiontitle{4. Порядок вычислений}
\begin{itemize}
\item Сначала выполняем действия в скобках.
\item Затем степени; далее умножение и деление; после них сложение и вычитание.
\item На одном уровне действуем слева направо, если нет причины удобнее сгруппировать слагаемые.
\end{itemize}

\sectiontitle{5. Текстовая задача: математическая модель}
\keyidea{Сначала определите величины и связи, затем составьте таблицу или выражение. Вычисления --- завершающий этап решения.}

\rulelabel{Количество и величина одной единицы} Если каждый кубический сантиметр вещества имеет массу \(q\) граммов, то при объёме \(V\) кубических сантиметров
\[
m=qV.
\]
Например, если \(q=11\)~г/см\(^3\), \(V=3\)~см\(^3\), то \(m=11\cdot3=33\)~г.

\tip{Умножать следует величину одной единицы на \emph{количество единиц}. Само по себе увеличение двух величин ещё не доказывает, что их нужно перемножать.}

\sectiontitle{6. Движение навстречу: чертёж и таблица}
При равномерном движении:
\[
S=vt,\qquad v=\frac St,\qquad t=\frac Sv\quad (v>0).
\]
\begin{center}
\begin{tikzpicture}[scale=0.92,>=Latex,font=\small]
\coordinate (A) at (0,0); \coordinate (M) at (6.0,0); \coordinate (B) at (12,0);
\draw[gray!65,line width=1pt] (A)--(B);
\fill[darkaccent] (A) circle (2.3pt) (B) circle (2.3pt);
\fill[accent] (M) circle (2.6pt);
\draw[accent,->,line width=1.4pt] (0.5,0.48)--(5.5,0.48);
\draw[accent,->,line width=1.4pt] (11.5,0.48)--(6.5,0.48);
\node[below=4pt] at (A) {Город A};
\node[below=4pt] at (M) {Встреча};
\node[below=4pt] at (B) {Город B};
\node[above=15pt] at (3,0) {\(S_1=v_1t\)};
\node[above=15pt] at (9,0) {\(S_2=v_2t\)};
\end{tikzpicture}
\end{center}
Два поезда одновременно выехали навстречу друг другу и встретились через \(3{,}35\)~ч.
Скорость первого --- 75~км/ч, второго --- \(\frac9{10}\) скорости первого.

\begin{center}
\renewcommand{\arraystretch}{1.4}
\begin{tabularx}{\linewidth}{@{}l>{\centering\arraybackslash}X>{\centering\arraybackslash}X>{\centering\arraybackslash}X@{}}
\toprule
Поезд & \(v\), км/ч & \(t\), ч & \(S\), км\\
\midrule
Первый & \(75\) & \(3{,}35\) & \(75\cdot3{,}35\)\\
Второй & \(\frac9{10}\cdot75=67{,}5\) & \(3{,}35\) & \(67{,}5\cdot3{,}35\)\\
\bottomrule
\end{tabularx}
\end{center}
\[
S_{\text{общ}}=S_1+S_2=(75+67{,}5)\cdot3{,}35
=\boxed{477{,}375\ \text{км}}.
\]
\tip{Если оба начали движение одновременно и встретились, время в пути у них одинаковое. Сумма пройденных путей равна начальному расстоянию.}

\sectiontitle{7. Если время движения разное}
Принцип сложения путей сохраняется, однако для каждого участника используем \emph{его собственное время}:
\[
S_{\text{общ}}=v_1t_1+v_2t_2.
\]
\rulelabel{Пример} Друзья двигались навстречу: первый \(1{,}4\)~ч со скоростью \(5{,}34\)~км/ч,
второй \(1{,}15\)~ч со скоростью \(6{,}1\)~км/ч.
\[
S=5{,}34\cdot1{,}4+6{,}1\cdot1{,}15
=7{,}476+7{,}015
=\boxed{14{,}491\ \text{км}}.
\]

\sectiontitle{8. Как найти дробь от числа}
\keyidea{Предлог «от» означает умножение:}
\[
\text{дробь от числа}=\text{дробь}\cdot\text{число}.
\]
\[
\frac27\ \text{от}\ 14=\frac27\cdot14=\boxed4,
\qquad
0{,}7\ \text{от}\ 40=0{,}7\cdot40=\boxed{28}.
\]
Если в одном выражении десятичная и смешанная дроби, удобно перевести их в обыкновенные:
\[
0{,}7\ \text{от}\ \mixed{1}{1}{2}
=\frac7{10}\cdot\frac32
=\boxed{\frac{21}{20}=\mixed{1}{1}{20}}.
\]

\sectiontitle{9. Как найти процент от числа}
Один процент --- это \(\frac1{100}\) целого. Следовательно,
\[
p\%\ \text{от}\ N=\frac p{100}\cdot N.
\]
\[
30\%\ \text{от}\ 40=\frac{30}{100}\cdot40=\boxed{12}.
\]
\[
55\%\ \text{от}\ 13{,}8
=\frac{55}{100}\cdot\frac{138}{10}
=\frac{759}{100}=\boxed{7{,}59}.
\]
\tip{Проценты сначала делите на 100, потом умножайте на исходное число.}

\sectiontitle{10. Три формы одной и той же доли}
\[
\boxed{\ 0{,}6=\frac35=60\%\ }.
\]
Поэтому расход \(0{,}6\) от 130~л, \(\frac35\) от 130~л и \(60\%\) от 130~л одинаков:
\[
0{,}6\cdot130
=\frac35\cdot130
=\frac{60}{100}\cdot130
=\boxed{78\ \text{л}}.
\]
\keyidea{Перевод десятичной дроби в обыкновенную или проценты меняет запись, но сохраняет величину доли.}

\sectiontitle{11. Распределительное свойство умножения}
Если перед скобкой стоит множитель \(k\), его можно умножить на \emph{каждый} член суммы или разности:
\[
(a+b)k=ak+bk,\qquad (a-b)k=ak-bk.
\]
Особенно удобно раскрыть скобки, когда внешнее число делится на знаменатели дробей:
\[
\begin{aligned}
\left(\frac57-\frac25\right)\cdot35
&=\frac57\cdot35-\frac25\cdot35\\
&=25-14=\boxed{11}.
\end{aligned}
\]
\[
\left(\frac49+\frac7{36}\right)\cdot36
=16+7=\boxed{23}.
\]
\tip{Если скобка содержит разность, при раскрытии знак минус сохраняется. Умножить необходимо каждый член скобки.}

\clearpage
\sectiontitle{Алгоритм: как составить модель текстовой задачи}
\vspace{4mm}
\begin{center}
\begin{tikzpicture}[
node distance=11mm and 11mm,
>={Latex[length=2.4mm]},
connector/.style={-Latex,draw=darkaccent,line width=0.85pt},
action/.style={draw=accent,rounded corners=2mm,align=center,text width=100mm,minimum height=13mm,inner sep=3mm},
choice/.style={diamond,draw=accent,aspect=2.5,align=center,text width=37mm,inner sep=1mm},
branch/.style={draw=accent,rounded corners=2mm,align=center,text width=52mm,minimum height=22mm,inner sep=3mm},
startend/.style={draw=darkaccent,ellipse,align=center,minimum width=50mm,minimum height=11mm},
every node/.style={font=\small}]
\node[startend] (st) {Прочитайте условие};
\node[action,below=of st] (data) {Выпишите величины, единицы измерения,\\известные и искомые данные};
\node[choice,below=15mm of data] (kind) {Задача на движение?};
\node[branch,below left=15mm and 5mm of kind] (move) {Да: составьте таблицу\\ «скорость --- время --- путь»;\\ \(S=vt\)};
\node[branch,below right=15mm and 5mm of kind] (other) {Нет: определите долю,\\ процент или величину\\ одной единицы};
\node[action,below=55mm of kind] (relation) {Запишите связь величин: сумму путей,\\ дробь от числа или произведение};
\node[action,below=of relation] (calc) {Вычислите, сохранив правильные единицы};
\node[action,below=of calc] (verify) {Проверьте ответ по условию\\ и оцените его разумность};
\node[startend,below=of verify] (end) {Запишите ответ};
\draw[connector] (st)--(data);
\draw[connector] (data)--(kind);
\draw[connector] (kind.south west) -| node[pos=0.30,above left]{\scriptsize да} (move.north);
\draw[connector] (kind.south east) -| node[pos=0.30,above right]{\scriptsize нет} (other.north);
\coordinate (join) at ([yshift=7mm]relation.north);
\draw[darkaccent,line width=0.8pt] (move.south) |- (join);
\draw[darkaccent,line width=0.8pt] (other.south) |- (join);
\draw[connector] (join)--(relation.north);
\draw[connector] (relation)--(calc);
\draw[connector] (calc)--(verify);
\draw[connector] (verify)--(end);
\end{tikzpicture}
\end{center}

\clearpage
\sectiontitle{Тренировочный блок · Путь А}
\textbf{10 заданий.} Выполняйте задания по порядку. В текстовых задачах записывайте используемую математическую модель. Решения вынесены из этого блока.

\begin{enumerate}
\item Вычислите, предварительно сократив дроби:
\[
\frac29\cdot\frac{27}{8}.
\]
\item Вычислите произведение трёх дробей:
\[
\frac3{10}\cdot\frac{25}{9}\cdot\frac65.
\]
\item Найдите значения двух выражений и сравните их:
\[
\left(\frac23\right)^2
\qquad\text{и}\qquad
\frac{2^2}{3}.
\]
\item Переведите смешанное число в неправильную дробь и вычислите:
\[
\mixed{2}{1}{3}\cdot\frac9{14}.
\]
\item Вычислите удобным способом, используя распределительное свойство:
\[
\left(\frac47+\frac3{14}\right)\cdot42.
\]
\item Найдите \(\frac5{12}\) от 96~л.
\item Найдите \(35\%\) от 240~кг.
\item Объём детали равен \(2{,}5\)~см\(^3\), а масса \(1\)~см\(^3\) материала --- \(11{,}2\)~г.
Найдите массу детали.
\item Двое друзей двигались навстречу друг другу и встретились. Первый шёл \(1{,}5\)~ч
со скоростью \(4{,}8\)~км/ч, второй --- \(1{,}2\)~ч со скоростью \(5{,}5\)~км/ч.
Каково было начальное расстояние между ними?
\item \textbf{Повышенная сложность.} Два поезда одновременно выехали навстречу друг другу
и встретились через \(2{,}4\)~ч. Скорость первого \(68\)~км/ч,
а скорость второго составляет \(\frac34\) скорости первого.
Найдите расстояние между городами.
\end{enumerate}

\clearpage
\sectiontitle{Ответы к тренировочному блоку}
\begin{center}
\renewcommand{\arraystretch}{1.68}
\begin{tabularx}{\linewidth}{@{}>{\centering\arraybackslash}p{15mm}X@{}}
\toprule
№ & Ответ\\
\midrule
1 & \(\displaystyle \frac34\).\\
2 & \(1\).\\
3 & \(\displaystyle \frac49\) и \(\displaystyle\frac43=\mixed{1}{1}{3}\);
второе выражение больше.\\
4 & \(\displaystyle \mixed{1}{1}{2}\).\\
5 & \(33\).\\
6 & \(40\)~л.\\
7 & \(84\)~кг.\\
8 & \(28\)~г.\\
9 & \(13{,}8\)~км.\\
10 & \(285{,}6\)~км.\\
\bottomrule
\end{tabularx}
\end{center}

\sectiontitle{Самопроверка}
Отметьте навыки, которыми уже уверенно владеете.
\begin{itemize}
\item[\(\square\)] Я сокращаю множители перед умножением дробей.
\item[\(\square\)] Я перевожу смешанное число в неправильную дробь.
\item[\(\square\)] Я различаю \(\left(\frac ab\right)^2\) и \(\frac{a^2}{b}\).
\item[\(\square\)] Я составляю таблицу «скорость --- время --- путь».
\item[\(\square\)] Я нахожу дробь и процент от числа умножением.
\item[\(\square\)] Я знаю, что \(0{,}6=\frac35=60\%\).
\item[\(\square\)] Я использую распределительное свойство для упрощения вычислений.
\end{itemize}
\end{document}
